
% !TEX root = main.tex

\subparagraph*{Shifting}

Given a grade $\sgr$ and a pair of grades $\ple{\Rgr,\rgr}$, where the first one is regarded as an assumption grade, we define the \emph{shif} of $\ple{\Rgr,\rgr}$ by $\sgr$ as the operation 
\begin{quote}
$\sgr\rsh\ple{\Rgr,\rgr} = \ple{\Rgr\rsum\sgr\rmul\rgr,\rgr}$. 
\end{quote}

\PGComm{Da qui in poi \`e tutto rifatto}

This operation extends to types and contexts as follows:
\begin{itemize}
\item $\sgr\rsh\Graded{\tau}{\Rgr}{\rgr} = \Graded{\tau}{\Rgr\rsum\sgr\rmul\rgr}{\rgr}$ ,
\item $\sgr\rsh (\CtxElem{\loc}{\Rgr}{\rgr}{\tau},\Sigma)= \CtxElem{\loc}{\Rgr\rsum\sgr\rmul\rgr}{\rgr}{\tau},(\sgr\rsh\Sigma)$ and $\sgr\rsh\emptyset=\emptyset$.
\end{itemize}
It is easy to show that shifting is monotone over types and contexts and 
\begin{quote}
$(\ast)\quad\sgr \rsh (\rgr\rmul\Sigma) = \rgr\rmul ((\sgr\rmul\rgr)\rsh\Sigma)$ and
\end{quote}

\begin{proposition}\label{prop:rshift}
If $\IsWFField\Sigma\mem\val\T$, then, for every grade $\sgr$ we have 
$\IsWFField{\sgr\rsh\Sigma}{\mem}{\val}{\sgr\rsh\T}$. 
\end{proposition}
\begin{proof}
The proof is by coinduction. 
The only non-trivial case is when $\val = \loc$. 
Consider a judgement 
$\IsWFField{\sgr\rsh\Sigma}{\mem}{\loc}{\Graded\C{\Rgr\rsum\sgr\rmul\rgr}{\rgr}}$ such that $\IsWFField\Sigma\mem\loc{\Graded\C\Rgr\rgr}$ is derivable. 
We have to show that this is the consequence of an instance of the rule \refToRule{t-ref-in-mem} where the premises are appropriately shifted. 
Consider a judgement 
$\IsWFField{\sgr\rsh\Sigma}{\mem}{\loc}{\Graded\C{\Rgr\rsum\sgr\rmul\rgr}{\rgr}}$ such that $\IsWFField\Sigma\mem\loc{\Graded\C\Rgr\rgr}$ is derivable. 
We have to show that this is the consequence of an instance of the rule \refToRule{t-ref-in-mem} where the premises are appropriately shifted. 
By hypothesis we know that 
$\mem(\loc) = \Pair{\_}
{\Obj\C{\GradedValInd{\val}{1}{\rgr},\ldots,\GradedValInd\val{n}{\rgr}}} $ where
$\fields{\C} = \Field{\Graded{\tau_1}{\Rgr_1}{\rgr_1}}{\f_1} \ldots \Field{\Graded{\tau_n}{\Rgr_n}{\rgr_n}}{\f_n}$
and 
\begin{enumerate}
\item \label{it1} $\CtxElem{\loc}{\Rgr}{\rgr}{\C}\ctxord \Sigma$
\end{enumerate} and,
for all $i \in 1..n$ we have $\IsWFField{\Sigma_i}{\mem}{\val_i}{\T_i}$ where
\begin{enumerate}\setcounter{enumi}{1}
%\item  \label{it2}$\IsWFField{\Sigma_i}{\mem}{\val_i}{\T_i}$
\item  \label{it3} $\T_i\subt\Graded{\tau_i}{\Rgr_i\ctxsum\Rgr\rmul\rgr_i}{\rgr_i}$
\item  \label{it4}$\rgr\rmul\Sigma_i\ctxord \Sigma$
\end{enumerate} 
Consider as premises the judgements 
$\IsWFField{\Sigma'_i}{\mem}{\val_i}{\T'_i}$ where $\Sigma'_i=(\sgr\rmul\rgr)\rsh\Sigma_i$ and $\T'_i=(\sgr\rmul\rgr)\rsh\T_i\ (i \in 1..n)$. For all $i \in 1..n$ we get 
\begin{enumerate}[(a)]
\item \label{iit1} $\rgr\rmul\Sigma'_i\ctxord\sgr\rsh\Sigma$ from (\ref{it4}), monotonicity of $\rsh$ and $(\ast)$ 
\item \label{iit2} $\T'_i\subt\Graded{\tau_i}{\Rgr_i\rsum\Rgr\rmul\rgr_i\rsum(\sgr\rmul\rgr\rmul\rgr_i)}{\rgr_i}=\Graded{\tau_i}{\Rgr_i\rsum\Rgr\rmul\rgr_i\rsum(\sgr\rmul\rgr\rmul\rgr_i)}{\rgr_i}= \Graded{\tau_i}{\Rgr_i\rsum (\Rgr\rsum\sgr\rmul\rgr)\rmul\rgr_i}{\rgr_i}$ from (\ref{it3}) and monotonicity of $\rsh$ 
\end{enumerate}
Moreover,
\begin{enumerate}[(a)]\setcounter{enumi}{2}
\item \label{iit3} $\sgr\rsh\CtxElem{\loc}{\Rgr}{\rgr}{\C}\ctxord \sgr\rsh\Sigma$ from  (\ref{it1}) and monotonicity of $\rsh$.
\end{enumerate}
From (\ref{iit1}), (\ref{iit2}) and (\ref{iit3}) we get the thesis.
\end{proof}

\begin{proposition}\label{prop:split-ty} 
If  $\IsWFField\Sigma\mem\val{\Graded\tau\Rgr{\rgr}}$  and  $\rgr_1\rsum\rgr_2\rord\rgr$ then 
there are $\Sigma_1$ and $\Sigma_2$ such that 
$\Sigma_1\ctxsum\Sigma_2\ctxord\Sigma$, 
$\IsWFField{\Sigma_1}{\mem}{\val}{\Graded\tau{\Rgr\rsum\rgr_2}{\rgr_1}}$ and 
$\IsWFField{\Sigma_2}{\mem}{\val}{\Graded\tau{\Rgr\rsum\rgr_1}{\rgr_2}}$. 
\end{proposition} 
\begin{proof}
The only non-trivial case is when $\val = \loc$. 
By \refToRule{t-ref-in-mem} we know that 
$\tau = \C$,  $\mem(\loc) = \Pair{\_}
{\Obj\C{\GradedValInd{\val}{1}{\sgr},\ldots,\GradedValInd\val{n}{\sgr}}} $ where
$\fields{\C} = \Field{\Graded{\tau_1}{\Sgr_1}{\sgr_1}}{\f_1} \ldots \Field{\Graded{\tau_n}{\Sgr_n}{\sgr_n}}{\f_n}$
and letting $\Sigma=\CtxElem{\loc}{\Tgr}{\tgr}{\C},\Sigma'$
\begin{enumerate}
\item \label{i1} $\CtxElem{\loc}{\Rgr}{\rgr}{\C}\ctxord \CtxElem{\loc}{\Tgr}{\tgr}{\C}$
\end{enumerate} and,
for all $i \in 1..n$ we have $\IsWFField{\Sigma'_i}{\mem}{\val_i}{\T_i}$ where
\begin{enumerate}\setcounter{enumi}{1}
%\item  \label{it2}$\IsWFField{\Sigma_i}{\mem}{\val_i}{\T_i}$
\item  \label{i2} $\T_i\subt\Graded{\tau_i}{\Sgr_i\ctxsum\Rgr\rmul\sgr_i}{\sgr_i}$
\item  \label{i3}$\rgr\rmul\Sigma'_i\ctxord \Sigma$
\end{enumerate} 
Let $\Sigma_1 = \rgr_2\rsh\Sigma$ and  $\Sigma''_i=\rgr_2\rsh\Sigma'_i$.
From  \cref{prop:rshift},  $\IsWFField{\Sigma''_i}{\mem}{\val_i}{\rgr_2\rsh\T_i}$ for all $i \in 1..n$.
Let $\Sigma=\CtxElem{\loc}{\Tgr}{\tgr}{\C},\Sigma'$
\begin{enumerate}[(a)]
\item \label{ii1} $\CtxElem{\loc}{\Rgr\rsum\rgr_2}{\rgr_1}{\C}\bpa\ctxord?\epa\CtxElem{\loc}{\Tgr\rsum\rgr_2\rmul\tgr}{\tgr}{\C}$ from (\ref{i1}) we get $\rgr_1\rord\tgr$ and $\Tgr\rord\Rgr$ \PGComm{but this is not enough!  Come si fa?}
\end{enumerate} 
and for all $i \in 1..n$ we have 
\begin{enumerate}[(a)]\setcounter{enumi}{1}
\item  \label{ii2} $\rgr_2\rsh\T_i\subt\rgr_2\rsh\Graded{\tau_i}{\Sgr_i\ctxsum\Rgr\rmul\sgr_i}{\sgr_i}=\Graded{\tau_i}{\Sgr_i\rsum\Rgr\rmul\sgr_i\rsum\rgr_2\rmul\sgr_i}{\sgr_i}=\Graded{\tau_i}{\Sgr_i\rsum(\Rgr\rsum\rgr_2)\rmul\sgr_i}{\sgr_i}$ from (\ref{i2}), monotonicity of $\rsh$ 
\item  \label{ii3}  $\rgr_1\ctxmul(\rgr_2\rsh\Sigma'_i)\bpa\ctxord?\epa \rgr_1\ctxmul((\rgr_2\rmul\rgr)\rsh\Sigma'_i)\rord\rgr\ctxmul((\rgr_2\rmul\rgr)\rsh\Sigma'_i)=\rgr_2\rsh(\rgr\ctxmul\Sigma'_i)\ctxord\rgr_2\rsh\Sigma$\\
\bpa assuming the property! $\rgr_2\bpa\ctxord?\epa\rgr_2\rmul\rgr$ when $\rgr_2\rord\rgr$ \epa
\end{enumerate} 
Finally to prove that $\Sigma_1\ctxsum\Sigma_2 \ctxord \Sigma$, observe that for all $\CtxElem{\loc'}{\Rgr'}{\rgr'}{\C'}\in\Sigma$, \\
\centerline{$\CtxElem{\loc'}{\Rgr'\rsum\rgr_2\rmul\rgr'}{\rgr'}{\C'}\ctxsum\CtxElem{\loc'}{\Rgr'\rsum\rgr_1\rmul\rgr'}{\rgr'}{\C'}$}
\PGComm{Che va bene per le assunzioni ma non per gli usi. Quindi questi $\Sigma_1$ e $\Sigma_2$ non vanno bene! }

Assuming the following rule
\[
\bpa
\NamedRule{t-ref-in-mem}{
  \IsWFField{\Sigma'_i}{\mem}{\val_i}{\T_i}\Space\forall i \in 1..n
}{ \IsWFField{ \rgr\rmul\Sigma}{\mem}{\loc}{\Graded{\C}{\Rgr}{\rgr}} }
{ \mem(\iota) =\Pair{\_}{\Obj{\C}{\GradedValInd{\val}{1}{\rgr},\ldots,\GradedValInd{\val}{n}{\rgr}}}\\
  \fields{\C}=\Field{\Graded{\tau_1}{\Sgr_1}{\sgr_1}}{\f_1} \ldots \Field{\Graded{\tau_n}{\Sgr_n}{\sgr_n}}{\f_n}\\
  \T_i\subt\Graded{\tau_i}{\Sgr_i\rsum\Rgr\rmul\sgr_i}{\sgr_i}\ \ \forall i\in 1..n\\
  (\ast)\ \CtxElem{\loc}{\Rgr}{\rgr}{\C}\ctxord\rgr\rmul\Sigma\\
  \Sigma'_i\ctxord\Sigma\ \ \forall i\in 1..n 
}
\epa
\]
with $\Sigma_1=\rgr_2\rsh\Sigma$ and $\Sigma_2=\rgr_1\rsh\Sigma$, let $\Sigma=\CtxElem{\loc}{\Tgr}{\tgr}{\C},\Sigma'$
 the  problem is to show $ (\ast)\ $, i.e.,
\[
\CtxElem{\loc}{\Rgr\rsum\rgr_2}{\rgr_1}{\C}\ctxord\CtxElem{\loc}{\Tgr\rsum\rgr_2\rmul\tgr}{\rgr_1\rmul\tgr}{\C}
\]
from $\rgr_1\rsum\rgr_2\rord\rgr$ and $\rgr\rord\rgr\rmul\tgr$ and $\Tgr\rord\Rgr$!
%
%
%Let us set 
%$\Sigma_1 = \CtxElem\loc{\Rgr\rsum\rgr_2}{\rgr_1}\C \ctxsum \rgr_1\rmul (\rgr_2\rsh\Sigma')$, then, by rule \refToRule{t-ref-in-mem}, we obtain 
%$\IsWFField{\Sigma_1}{\mem}{\loc}{\Graded\C{\Rgr\rsum\rgr_2}{\rgr_1}}$. 
%In a similar way we derive $\IsWFField{\Sigma_2}{\mem}{\loc}{\Graded\C{\Rgr\rsum\rgr_1}{\rgr_2}}$ where $\Sigma_2  = \CtxElem\loc{\Rgr\rsum\rgr_1}{\rgr_2}\C \ctxsum \rgr_2\rmul(\rgr_1\rsh\Sigma')$. 
%Therefore, to conclude, we just have to prove that $\Sigma_1\ctxsum\Sigma_2 \ctxord \Sigma$, but this is immediate by construction, as in each side we have increased the assumption grade of each reference exactly by its usage in the other side and  $\rgr_1\rsum\rgr_2\rord\rgr$. 
%


\end{proof}


